\chapter{Background and Related Work}
\label{chap:literature}

\noindent
This chapter surveys the scientific and engineering foundations on which
AnsibleGuard is built. Section~\ref{sec:iac_bg} introduces Infrastructure as
Code with emphasis on Ansible's execution model. Section~\ref{sec:smells_bg}
covers code smells and their extension to \ac{IaC}. Section~\ref{sec:checklist_bg}
presents the OpenStack code-review checklist paper that directly motivates
AnsibleGuard's rule taxonomy. Section~\ref{sec:static_bg} surveys static analysis
tools for Ansible. Section~\ref{sec:empirical_bg} reviews empirical \ac{IaC} quality
studies. Section~\ref{sec:gaps_bg} synthesizes the research gaps.

% ---------------------------------------------------------------
\section{Infrastructure as Code}
\label{sec:iac_bg}
% ---------------------------------------------------------------

\subsection{Definition and Historical Context}

\textit{Infrastructure as Code} is the practice of expressing every
configuration and provisioning decision in machine-readable files that are
version-controlled, reviewed, tested, and deployed through the same automation
pipelines used for application software \cite{humble2010continuous}. The term
emerged from the DevOps movement of the late 2000s, first articulated by Patrick
Debois at the Agile 2008 conference \cite{debois2008agile}, as teams began
applying software-engineering discipline to the manual, error-prone work of
server configuration. Before \ac{IaC}, a production environment could drift from
its intended state silently; after a misconfigured manual change, reproducing the
environment from scratch was a multi-day effort. \ac{IaC} addresses both problems:
all intended state is in version control, and a clean environment can be
re-provisioned from that state at any time.

\medskip\noindent
The practical consequences are significant. Large cloud deployments —
involving hundreds of machines, dozens of services, and multiple geographic
regions — are manageable only when their configuration is code. A pull request
on an \ac{IaC} repository is not a documentation change; it is a production
deployment proposal, and it must be reviewed with the same rigor as application
code \cite{bessghaier2025iac}.

\subsection{The IaC Tool Landscape}

Several tools have emerged with distinct approaches to expressing and enforcing
infrastructure state:

\begin{itemize}
  \item \textbf{Ansible} (Red Hat) — agentless, push-based, YAML playbooks
        describing operations to execute. The control machine connects over
        \ac{SSH}; no resident daemon is required on managed nodes.
  \item \textbf{Puppet} (Perforce) — agent-based, declarative \ac{DSL}. Resources
        declare a desired state that Puppet's agent enforces on a periodic schedule.
  \item \textbf{Chef} (Progress) — agent-based Ruby \ac{DSL}; procedural recipes
        and declarative resources for operators with Ruby backgrounds.
  \item \textbf{Terraform} (HashiCorp) — declarative HCL focused on cloud resource
        provisioning. Explicit plan/apply with state files.
  \item \textbf{SaltStack} (VMware) — event-driven, can run agent-based (minion)
        or agentless. Stronger real-time monitoring than Ansible.
\end{itemize}

\noindent
Puppet has been the most studied in the empirical quality literature
\cite{sharma2016puppet,bessghaier2024saner}. Ansible, despite leading industry
adoption surveys, has received comparatively less formal study — a gap that
motivates this work.

\subsection{Ansible: Architecture and Execution Model}
\label{sec:ansible_arch}

Ansible's core execution model is simpler than that of agent-based tools but
places more responsibility on the playbook author for idempotency. Its key
constructs are:

\begin{description}
  \item[Module.] An atomic operation encapsulated in a Python plugin. Ansible
    ships over 3,000 built-in modules in the \code{ansible.builtin} collection.
    Community collections (e.g., \code{community.general}, \code{community.aws})
    add thousands more. Modules are identified by their \textit{Fully Qualified
    Collection Name} (\ac{FQCN}): \code{ansible.builtin.copy},
    \code{community.aws.s3\_bucket}. Short names (\code{copy}, \code{s3\_bucket})
    are deprecated since Ansible 2.10.

  \item[Task.] A mapping in \ac{YAML} that invokes one module with arguments.
    A task may carry metadata: \code{name}, \code{when} conditions,
    \code{register} (capture output), \code{notify} (fire a handler),
    \code{tags}, \code{become} (privilege escalation), \code{loop},
    \code{ignore\_errors}, \code{no\_log}, and others.

  \item[Play.] A mapping that binds a set of tasks to a target host pattern.
    A play defines remote user, privilege-escalation settings, variable
    overrides, and the ordered list of \code{pre\_tasks}, \code{roles},
    \code{tasks}, and \code{post\_tasks}.

  \item[Playbook.] A YAML file containing one or more plays. Operators execute
    playbooks with the \code{ansible-playbook} command.

  \item[Role.] A self-contained reusable component with a standardized
    directory layout (Table~\ref{tab:role_layout}). Roles are the unit of
    sharing on Ansible Galaxy\cite{ansible_galaxy}, which hosts over 36,000
    community roles.

  \item[Handler.] A task that executes at the end of a play only when notified
    by at least one task. Used for service restarts, cache flushes, or other
    post-change reactions.

  \item[Variable.] Ansible defines 22 distinct variable-precedence levels,
    from role defaults (lowest) through command-line \code{--extra-vars}
    (highest). This complexity is the source of several AnsibleGuard rules:
    \code{role-variable-precedence-collision}, \code{unguarded-variable}, and
    \code{variable-use-before-definition}.
\end{description}

\begin{table}[htbp]
\centering
\caption{Standard Ansible role directory layout.}
\label{tab:role_layout}
\small
\begin{tabular}{ll}
\toprule
\textbf{Directory / File} & \textbf{Purpose} \\
\midrule
\code{tasks/main.yml}       & Default task entry point \\
\code{handlers/main.yml}    & Handler definitions \\
\code{defaults/main.yml}    & Role defaults (lowest precedence) \\
\code{vars/main.yml}        & Role variables (overrides play/task) \\
\code{templates/}           & Jinja2 template files \\
\code{files/}               & Static files to copy \\
\code{meta/main.yml}        & Role metadata, dependencies, Galaxy info \\
\code{tests/}               & Molecule or other test scenarios \\
\code{README.md}            & Required for Galaxy publication \\
\bottomrule
\end{tabular}
\end{table}

\subsection{Why Ansible Code Is Hard to Analyze}

Ansible's simplicity is also its vulnerability to quality problems. Because
any valid YAML file that loosely follows the expected schema is accepted as a
playbook, there are no compile-time or type-system gates. Key analysis
challenges include:

\begin{itemize}
  \item \textbf{No type system.} Variables are untyped; their value at runtime
        depends on which of 22 precedence levels wins, including sources
        invisible at scan time (inventory, \code{--extra-vars}, facts, vault).
  \item \textbf{Jinja2 templating inside YAML.} Expressions like
        \code{\{\{ my\_var | default('fallback') \}\}} must be parsed as
        Jinja2 within a YAML string, requiring a two-level parser.
  \item \textbf{Dynamic includes.} \code{include\_tasks} with a Jinja2 path
        argument cannot be resolved statically; the included file's tasks
        are invisible to static analysis.
  \item \textbf{Idempotency is the author's responsibility.} Unlike Puppet,
        Ansible does not guarantee idempotency; whether a
        \code{command:} task changes the system depends on the author's use
        of \code{creates:}, \code{removes:}, or \code{changed\_when:}.
\end{itemize}

% ---------------------------------------------------------------
\section{Code Smells and IaC Quality}
\label{sec:smells_bg}
% ---------------------------------------------------------------

\subsection{The Code Smell Concept}

The term \textit{code smell} was coined by Kent Beck and popularized by Martin
Fowler \cite{fowler1999refactoring}. A smell is not a bug — the code runs —
but its structure makes it harder to understand, evolve, or maintain than it
needs to be. Fowler's catalogue pairs each smell with a named refactoring.
Steve McConnell \cite{mcconnell2004codecomplete} frames the same intuition in
terms of \textit{code quality attributes}: correctness, readability, robustness,
and maintainability. Both works established that structural quality is
measurable, documentable, and worth automating.

\subsection{IaC Smell Studies}

The adaptation of smells to \ac{IaC} began with Sharma, Fragkoulis, and
Spinellis \cite{sharma2016puppet}, who identified ten configuration smells in
Puppet manifests using an automated detector (Puppeteer). Their analysis of
2,214 Puppet repositories found that misplaced attributes, long resources, and
missing documentation occur frequently and correlate with subsequent
defect-fixing commits — providing the first empirical evidence that \ac{IaC}
smells matter in practice.

\medskip\noindent
Bessghaier et al.\ \cite{bessghaier2024saner} extended this line of work in
their SANER 2024 paper, studying the \textit{prevalence, co-occurrence, and
impact} of 12 Puppet smell types across 82 projects using 19,641 commits. Their
study demonstrated that smelly Puppet files are more change-prone and
defect-prone than clean files, lending empirical weight to the practical importance
of smell detection. The 12 smells comprise four implementation smells (complex
expression, incomplete tasks, deprecated statement, incomplete conditional) and
eight design smells (multifaceted abstraction, insufficient modularization,
tight coupling, weakened modularity, unstructured module, missing dependency,
dense structure, deficient encapsulation). AnsibleGuard directly maps these
12 smells to equivalent Ansible constructs.

\medskip\noindent
Rahman, Parnin, and Williams \cite{rahman2019security} identified seven security
smells in Ansible and Chef scripts, reporting precision between 72\,\% and
93\,\% per smell. Their protocol informs a future labelled evaluation of
AnsibleGuard, but those figures are not directly comparable to the unlabelled
occurrence distributions in Chapter~\ref{chap:results}.

\medskip\noindent
Opdebeeck, Zerouali, and De Roover \cite{opdebeeck2023ansible} studied 8,065
Ansible Galaxy roles and found that 97\,\% contained at least one smell from
a combined catalogue. Documentation and role-variable scoping problems dominate.
Kumari and Rahman \cite{kumari2021ansible} further studied maintainability smells
in Ansible specifically, cataloguing patterns that increase change cost without
providing correctness benefits.

% ---------------------------------------------------------------
\section{The OpenStack IaC Code-Review Checklist}
\label{sec:checklist_bg}
% ---------------------------------------------------------------

The most direct motivator for AnsibleGuard's rule taxonomy is the study by
Bessghaier et al.\ \cite{bessghaier2025iac}, published in \textit{Empirical
Software Engineering} 30:106, 2025. The study performed a large-scale empirical
analysis of Ansible code review on OpenStack projects:

\begin{itemize}
  \item \textbf{Dataset:} 300,000+ OpenStack Gerrit review comments and 379
        manually sampled merged \ac{IaC} changes.
  \item \textbf{Methodology:} Two authors independently coded the sample with
        Krippendorff's $\alpha = 0.99$ after discussion; nine of the 100 most
        active OpenStack Ansible reviewers validated the checklist (98\,\%
        positive ratings across 324 responses).
  \item \textbf{Output:} 38 review criteria grouped into seven topics —
        Documentation, Logic, Best Practices, Compatibility, Dependencies,
        Test, and Security.
\end{itemize}

\medskip\noindent
The checklist derived from this study is the ground truth for AnsibleGuard's
design. Table~\ref{tab:checklist_summary} summarizes the seven topics and their
coverage in AnsibleGuard (full mapping in Chapter~\ref{chap:implementation}).

\begin{table}[htbp]
\centering
\caption{Summary of the 38-criterion OpenStack IaC review checklist and AnsibleGuard coverage.}
\label{tab:checklist_summary}
\small
\begin{tabular}{lrrrr}
\toprule
\textbf{Topic} & \textbf{Criteria} & \textbf{Yes} & \textbf{Partial} & \textbf{No} \\
\midrule
Documentation    & 5  & 1 & 3 & 1 \\
Logic            & 11 & 5 & 5 & 1 \\
Best Practices   & 9  & 6 & 2 & 1 \\
Compatibility    & 4  & 0 & 4 & 0 \\
Dependencies     & 2  & 0 & 2 & 0 \\
Test             & 3  & 1 & 2 & 0 \\
Security         & 4  & 3 & 1 & 0 \\
\midrule
\textbf{Total}   & \textbf{38} & \textbf{15} & \textbf{18} & \textbf{5} \\
\bottomrule
\end{tabular}
\end{table}

\noindent
The five criteria with no automation are: blank line before a configuration
section (a source-format convention, not an Ansible semantic), setting
deprecated parameters to UNDEF (a Puppet-specific construct inapplicable to
Ansible), following OpenStack IRC conventions (a process rule, not a YAML
property), testing assigned values for correctness (requires test execution),
and parameter-value compatibility (requires service-specific knowledge). These
are permanently excluded from AnsibleGuard's scope.

% ---------------------------------------------------------------
\section{Static Analysis Tools for Ansible}
\label{sec:static_bg}
% ---------------------------------------------------------------

\subsection{Precision and Recall}

Static analysis trades off precision against recall. A tool that reports
every file as defective has 100\,\% recall but potentially near-0\,\% precision;
practitioners may ignore it. Precision therefore requires a labelled sample,
whereas recall requires known relevant defects. Neither can be inferred from a
raw finding count. This report treats both as future evaluation outcomes.

\subsection{Existing Ansible Quality Tools}

\textbf{ansible-lint}\cite{ansiblelint_docs} is the primary Ansible quality
tool. Maintained by the Ansible community (Red Hat), it implements
approximately 180 built-in rules covering \ac{FQCN} compliance, variable naming,
task structure, and syntax correctness. Its limitations relative to AnsibleGuard:
it has no size/complexity rules, no cross-file analysis (handler resolution,
variable scoping), no pull-request diff mode, and no empirical precision study.

\medskip\noindent
\textbf{yamllint}\cite{yamllint_docs} enforces YAML formatting (indentation,
trailing spaces, colon spacing) without any Ansible-specific semantic understanding.
AnsibleGuard integrates both tools as providers, normalizing their findings into
the same output schema as internal rules.

\medskip\noindent
\textbf{GitLeaks}\cite{gitleaks_github} detects secrets in repository trees
and Git history using regular expressions and entropy heuristics. It covers
hard-coded credentials, API keys, and other sensitive material. AnsibleGuard
integrates GitLeaks as an external provider, with a conservative noise filter
that suppresses findings matching known checksum or comment patterns.

\medskip\noindent
\textbf{SLIC} \cite{rahman2019security} is a research tool implementing the
seven security smells for Ansible and Chef. It provides the most directly
comparable published precision baseline (72–93\,\% per smell). Unlike
AnsibleGuard, SLIC covers only security and has not been validated at scale
with a dual-reviewer protocol.

\subsection{The SARIF Interchange Format}

SARIF 2.1.0 \cite{oasis2019sarif} is an OASIS standard JSON schema for static
analysis results. Its significance for CI/CD integration is that GitHub Code
Scanning\cite{github_sarif}, VS Code, and JetBrains all consume SARIF files
natively, rendering findings as inline pull-request line annotations without
custom post-processing. AnsibleGuard produces SARIF with:
\begin{itemize}
  \item repository-relative \code{artifactLocation.uri} (required by GitHub),
  \item stable \code{partialFingerprints.ansibleguardSemanticId/v1} (required
        for result identity after line movements), and
  \item \code{ruleId} mapped to the native ansible-lint rule (not the generic
        \code{ansible-lint-issue} envelope) so each lint check gets a distinct
        dismissal scope in GitHub Code Scanning.
\end{itemize}

% ---------------------------------------------------------------
\section{Empirical IaC Quality Studies}
\label{sec:empirical_bg}
% ---------------------------------------------------------------

\subsection{Defect Studies}

Rahman and Williams \cite{rahman2015iac} analyzed 1,726 defect-fixing commits
across Ansible and Chef repositories, finding that security and dependency
defects carry the highest fixing cost. Their follow-up \cite{rahman2020defective}
applied source-code metrics (change frequency, code churn) to predict defect-prone
files with AUC $> 0.75$.

\subsection{Precision Measurement Protocols}

The standard protocol for measuring static analysis precision involves:
(1) stratified sampling to ensure representation of all rules, (2) two
independent reviewers, (3) adjudication of disagreements by a third reviewer,
and (4) Wilson-score confidence intervals \cite{wilson1927probable} rather than
normal approximations, which are inaccurate for proportions near 0 or 1.
Cohen's $\kappa$ \cite{cohen1960kappa} quantifies inter-rater agreement. Rahman
et al.\ \cite{rahman2019security} used this protocol with $\kappa = 0.92$ for
their seven security smells. The present work specifies this protocol for a
future labelled study; it does not claim that independent review and
adjudication have already been completed.

% ---------------------------------------------------------------
\section{Research Gaps and Positioning}
\label{sec:gaps_bg}
% ---------------------------------------------------------------

\noindent
Three gaps motivate AnsibleGuard:

\begin{enumerate}
  \item \textbf{Breadth.} No single Ansible tool covers structural, security,
        correctness, documentation, and testing concerns in a unified pipeline.
  \item \textbf{Empirical rigor.} Tool behaviour, provider health, labelled
        accuracy, and ecosystem prevalence are often conflated; a reproducible
        evaluation must keep these claims separate.
  \item \textbf{CI integration.} No existing Ansible tool offers a pull-request
        diff mode with SARIF output and GitHub inline annotations out of the box.
\end{enumerate}

Table~\ref{tab:tool_comparison_bg} positions AnsibleGuard against the closest
related tools.

\begin{table}[htbp]
\centering
\caption{Feature comparison of Ansible quality tools.}
\label{tab:tool_comparison_bg}
\small
\begin{tabular}{lccccc}
\toprule
\textbf{Feature} & \rotatebox{70}{\textbf{ansible-lint}} & \rotatebox{70}{\textbf{yamllint}} & \rotatebox{70}{\textbf{GitLeaks}} & \rotatebox{70}{\textbf{SLIC}} & \rotatebox{70}{\textbf{AnsibleGuard}} \\
\midrule
Style / naming      & \checkmark & \checkmark &  &  & \checkmark \\
Correctness         & partial    &            &  & partial & \checkmark \\
Security            & partial    &            & \checkmark & \checkmark & \checkmark \\
Size / complexity   &            &            &  &  & \checkmark \\
Testing / checklist &            &            &  &  & \checkmark \\
Cross-file analysis &            &            &  &  & \checkmark \\
Symbol index        &            &            &  &  & \checkmark \\
SARIF output        & \checkmark &            &  &  & \checkmark \\
PR diff mode        &            &            &  &  & \checkmark \\
Precision study     &            &            &  & \checkmark & \checkmark \\
149-repo corpus     &            &            &  &  & \checkmark \\
\bottomrule
\end{tabular}
\end{table}
